%-----------------------------------
\section{\textit{Survey}}
\label{sec_tipo_survey}

\subsection{Caracterização}

O \textit{survey}, ou levantamento, coleta informações de um grupo de pessoas, geralmente por meio de questionário, para descrever características, opiniões ou comportamentos de uma população \cite{kitchenham2008surveys}. Exemplos em Sistemas de Informação incluem a adoção de tecnologias por pequenas empresas, a percepção de estudantes sobre ferramentas de ensino e o perfil de profissionais de desenvolvimento.

Os cuidados principais são a definição da população e da amostra, a construção e a validação do instrumento (por exemplo, com um pré-teste) e os aspectos éticos. Pesquisas com seres humanos exigem o Termo de Consentimento Livre e Esclarecido (TCLE) e, conforme o caso, a aprovação de um Comitê de Ética em Pesquisa (CEP) \cite{cns2016resolucao510}, além do tratamento adequado dos dados pessoais \cite{brasil2018lgpd}.

\subsection{Estrutura típica do trabalho}

\begin{quadro}[htb]
	\centering
	\caption{Estrutura típica de um TCC de \textit{survey}.}
	\label{qua_estrutura_survey}
	\small
	\begin{tabularx}{\textwidth}{|l|X|}
		\hline
		\textbf{Capítulo} & \textbf{Seções típicas} \\ \hline
		Metodologia & Questões de pesquisa; população e amostra; instrumento de coleta e escalas; validação do instrumento; procedimento de coleta; análise dos dados; aspectos éticos. \\ \hline
		Resultados e discussão & Perfil dos respondentes; resultados por dimensão investigada; respostas às questões de pesquisa; ameaças à validade. \\ \hline
		Apêndices e anexos & Questionário completo e TCLE (apêndices); parecer do CEP (anexo). \\ \hline
	\end{tabularx}
	\fonte{Elaborado pelo autor.}
\end{quadro}

O Quadro~\ref{qua_estrutura_survey} resume a estrutura desse tipo de trabalho. Modelos de questionário e de TCLE estão no Apêndice~\ref{apC_questionario}.

\subsection{Exemplo: respostas em escala Likert}

Questões em escala Likert são frequentemente apresentadas em gráfico de barras empilhadas, como na Figura~\ref{fig_likert}, que mostra a distribuição percentual das respostas a cada afirmação.

\begin{figure}[htb]
	\centering
	\caption{Percepção dos respondentes sobre o sistema acadêmico (dados fictícios, $n = 120$).}
	\label{fig_likert}
	\begin{tikzpicture}[x=0.09cm, y=0.8cm, font=\small]
		% Cores (tons de cinza para impressão)
		\def\cA{black!80}\def\cB{black!55}\def\cC{black!30}\def\cD{black!15}\def\cE{black!5}
		% \likertbar{linha}{rótulo}{DT}{D}{N}{C}{CT}
		\newcommand{\likertbar}[7]{%
			\node[anchor=east] at (-2,#1) {#2};
			\pgfmathsetmacro{\a}{#3}\pgfmathsetmacro{\b}{\a+#4}\pgfmathsetmacro{\c}{\b+#5}\pgfmathsetmacro{\d}{\c+#6}
			\fill[\cA] (0,#1-0.3) rectangle (\a,#1+0.3);
			\fill[\cB] (\a,#1-0.3) rectangle (\b,#1+0.3);
			\fill[\cC] (\b,#1-0.3) rectangle (\c,#1+0.3);
			\fill[\cD] (\c,#1-0.3) rectangle (\d,#1+0.3);
			\fill[\cE] (\d,#1-0.3) rectangle (100,#1+0.3);
			\draw (0,#1-0.3) rectangle (100,#1+0.3);
		}
		\likertbar{3}{É fácil de usar}{5}{10}{20}{40}{25}
		\likertbar{2}{É rápido}{10}{20}{30}{30}{10}
		\likertbar{1}{Atende às necessidades}{5}{15}{25}{35}{20}
		% Eixo
		\draw (0,0.5) -- (100,0.5);
		\foreach \x in {0,25,50,75,100} {\draw (\x,0.5) -- (\x,0.4) node[below] {\x\%};}
		% Legenda
		\foreach \cor/\rot [count=\i from 0] in {\cA/Discordo totalmente, \cB/Discordo, \cC/Neutro, \cD/Concordo, \cE/Concordo totalmente} {
			\pgfmathsetmacro{\lx}{mod(\i,3)*40-25}
			\pgfmathsetmacro{\ly}{-0.6-floor(\i/3)*0.5}
			\fill[\cor, draw=black] (\lx,\ly-0.12) rectangle (\lx+3,\ly+0.12);
			\node[anchor=west, font=\scriptsize] at (\lx+3.5,\ly) {\rot};
		}
	\end{tikzpicture}
	\fonte{Elaborado pelo autor.}
\end{figure}
